\documentclass[11pt]{article}

\usepackage[margin=1in]{geometry}
\usepackage{amsmath, amssymb, amsthm}
\usepackage{booktabs}
\usepackage{hyperref}
\usepackage{enumitem}
\usepackage{xcolor}
\usepackage[colorinlistoftodos]{todonotes}

\newcommand{\kQ}{\kappa_{Q}}
\newcommand{\Pij}{\Pi_{ij}}
\newcommand{\Sj}{\Sigma_{j}}
\newcommand{\Cj}{C_{j}}
\newcommand{\E}{\mathbb{E}}
\newcommand{\Var}{\operatorname{Var}}
\newcommand{\Cov}{\operatorname{Cov}}
\newcommand{\Rb}{\mathbb{R}}
\newcommand{\Nrm}{\mathcal{N}}
\newcommand{\nuhat}{\hat\nu}
\newcommand{\MAE}{\mathrm{MAE}}
\newcommand{\MAPE}{\mathrm{MAPE}}

\title{\textbf{Qualifier Exhaustion: Can the \(\kQ\)-framework
Accommodate the Heavy-Tail / Forecast-Accuracy Coupling?}\\[4pt]
\large A staged exhaustion memo: \(\kQ\)-only \(\to\) named extensions
\(\to\) coherent generalisations}
\author{MITACS dynamics project --- design memo, not a result artifact}
\date{2026-05-20}

\begin{document}
\maketitle

\begin{center}
\fbox{\parbox{0.92\linewidth}{\centering
\textbf{PROVENANCE-GRADE: INSPECTION-ONLY.} This memo is a framework
analysis; it issues no estimator change and no claim about Ontario
data beyond what is already established and provenanced in the
referenced result artifacts. Conclusions are statements about what
the \(\kQ\) framework does and does not license, not about what the
estimator produces.}}
\end{center}

\begin{abstract}
The investigation in \S4.2 of the application paper closed three
climatology-axis experiments (seasonal-vs-raw; scale-only @ elbow;
scale-only @ forced seasonal dims) on the same parametric object
\((\mu_{m,h}, \sigma_{m,h})\) and reached a HOLD verdict on
de-seasonalisation removal with branch (iii) confirmed (the additive
sub-step is doing real work independent of dimension). Two engineering
candidates remain inside that family --- additive seasonal with
exogenous scale (\(B\)); Fourier-decomposition scale layer (\(C\)) ---
but both feed transformed series into the same Gaussian local kernel,
so neither in principle reaches the heavy-tail cost; they can only
redistribute it. The question this memo answers is whether the
\(\kQ\) framework permits any method that accommodates the coupling,
and if not, why.

The exhaustion is \emph{staged}. At each stage we attempt to settle
every one of the three qualifiers --- \textbf{Register},
\textbf{Imposed Gaussianity}, \textbf{Single step} --- as
\textsc{open} (the relaxation lives inside the stage's licensing),
\textsc{foreclosed-at-stage} (the relaxation is incoherent at this
stage but may live at the next), or \textsc{requires-experiment}
(open in principle; needs a designed gate). We escalate only when all
three reach a definite verdict at the current stage. The three stages
are: (1) strict \(\kQ\)-only (no redefinition); (2) \(\kQ\) +
extensions the framework's own language admits (e.g.\ law-agnostic
local kernels); (3) any Markov-semigroup-coherent estimator class
(may leave the \(\kQ\) register).

The memo is paper-only --- experiments fall out of its gaps as
separate work. Integration into the application paper (sibling tex
vs.\ inline \S4.3) is deferred until the verdict shape is visible.
\end{abstract}

\section{Setup and inheritance}
\label{sec:setup}

\paragraph{What we inherit from the paper.} The application paper's
\S2.2 defines \(\kQ\) via Rokhlin disintegration of the joint law
\(Q\) over the coarsened \(\sigma\)-algebra induced by the delay
embedding, gives the Chapman--Kolmogorov-derived Markov semigroup
\(\{\Pi_t\}\), and states the Koopman--Perron duality. The estimator
\(\Pij\) targets the one-step kernel: per anchor \(j\), a local
linear drift \(\Cj\) and a Gaussian diffusion \(\Sj\) define
\(\Pij(\cdot \mid x) = \Nrm(x\Cj,\; \Sj)\). The three qualifiers
that \S2.2 documents as the ways \(\Pij\) departs from \(\kQ\) are:

\begin{description}[leftmargin=2em]
  \item[Register.] \(\kQ\) is a disclosure object (the conditional
  law that exists by Rokhlin); \(\Pij\) is a parametric construction
  inside that disclosure --- a specific local-linear-Gaussian
  realisation. The framework permits other realisations.
  \item[Imposed Gaussianity.] The conditional \emph{law} is fit as
  Gaussian by estimator design (WLS gives \(\Cj\), residual
  covariance gives \(\Sj\), and these together specify a Gaussian
  predictive law). \(\kQ\) itself is law-agnostic --- whatever the
  conditional distribution of \(Q\) is, that is what \(\kQ\)
  discloses. Gaussianity is not a \(\kQ\)-side commitment.
  \item[Single step.] \(\Pij\) targets \(\Pi_1\) only; multi-step
  predictions are obtained by composition (\(\Pi_1^k\)), which carries
  a propagation gate (Task~\#11/12) but does not refit per horizon.
  \(\kQ\) supports the full semigroup \(\{\Pi_t\}\); the estimator
  occupies only the one-step generator.
\end{description}

\paragraph{What we inherit from results.}
\begin{itemize}[leftmargin=2em]
  \item Rebaseline (\textsc{claim-grade}): one-step scalar
  innovation excess kurtosis \(\approx 24\)--\(33\) across day-types
  on the intra-day-conditioned residual; tail ratio \(\approx
  2.4\)--\(3.8\). This is the heavy-tail evidence the memo is trying
  to accommodate.
  \item Error-decomposition (\textsc{method/design-grade}): on the
  validated full population (132 days via \(d_{\max}\)-augmented
  propagation), \(M_1 \approx 4.8\%\), \(M_2 \approx 2.3\%\); CIs
  overlap; neither drift- nor diffusion-deficiency dominates.
  \item \S4.2 disambiguation (\textsc{inspection-only}):
  forecast-MAPE ranking at matched dim \(\{2,4,3\}\) is seasonal
  \(3.86\% \to\) scale-only \(13.13\% \to\) no-transform \(13.57\%\).
  The additive \(\mu_{m,h}\) step alone carries
  \(\approx 9.3\) pp of the \(\approx 9.7\) pp total cost; the
  scale step \(\approx 0.4\) pp; the dimension shift
  \(\approx 1.1\) pp. Branch (iii) confirmed.
  \item \#41 (\textsc{resolved}): no prediction-error
  cross-validation localises a bandwidth on this estimator on any
  tested system; the operator-as-run is a global linear drift. Any
  Stage-2/3 lever that adds a hyperparameter is subject to this
  failure mode if naively selected by CV.
\end{itemize}

\paragraph{The coupling the memo must accommodate.} The three
empirical facts taken together --- strongly non-Gaussian one-step
innovation; both drift and diffusion deficiencies materialise; the
additive level step is irreducibly entangled with forecast accuracy
in the parametric \((\mu, \sigma)\) family --- describe a regime in
which a Gaussian local kernel fed a level-and-scale-stripped
embedding is fundamentally a proxy. The question is whether the
framework permits any escape from the proxy regime.

\paragraph{Exit condition (locked).} Each stage produces a
three-row table: \{Register, Imposed Gaussianity, Single step\}
\(\times\) \{verdict, mechanism, evidence in hand, what would close
it\}. We escalate to the next stage only when all three rows hold a
definite verdict at the current stage.

\section{Stage 1 --- Strict \(\kQ\)-only (no redefinition)}
\label{sec:stage1}

\paragraph{Licensing (partition declaration).} The \(\kQ\)
disclosure is law-agnostic and supports the full semigroup
\(\{\Pi_t\}\); it does not by itself commit to any specific
estimator class. To make the staging informative rather than
tautological we declare a \emph{partition} at Stage 1: the local
estimator class stays \emph{local-linear-Gaussian} and the disclosure
register stays unchanged; what varies is (i) the \emph{state} the
realisation is built on (the embedding that defines the conditioning
\(\sigma\)-algebra) and (ii) the \emph{horizon} \(\Pi_t\) being
targeted. Stage 2 relaxes the estimator class explicitly; Stage 3
relaxes the framework. Stage 1's foreclosures of Imposed Gaussianity
that follow are foreclosures \emph{at this partition}, not
\(\kQ\)-side commitments.

\paragraph{Register at Stage 1.} \(\kQ\) is defined via the Rokhlin
disintegration of \(Q\) over the \(\sigma\)-algebra induced by the
observation pair \((Q, F)\). The embedding is what \emph{defines}
that \(\sigma\)-algebra: changing the embedding gives a different
\((Q, F)\) and therefore a different disclosure object
\(\kappa_{Q', F'} \neq \kappa_{Q, F}\). The framework permits this --- the
Rokhlin theorem applies to any well-defined observation pair --- but
the resulting estimand is a different conditional law, not a
re-reading of the same one. Adding a derivative coordinate, an
explicit phase coordinate, or an exogenous covariate is therefore a
substantive Stage-1 move: it stays inside the local-linear-Gaussian
class and inside the disclosure register but it targets
\(\kappa_{Q,\mathrm{new}}\). The §6 conclusion already names this as the
higher-leverage path. \textbf{Verdict: OPEN, responsive, with
estimand-change caveat.}

\paragraph{Imposed Gaussianity at Stage 1.} Under the partition, the
estimator class is held fixed at local-linear-Gaussian. Relaxing the
local kernel (heavier-tailed parametric form; non-parametric local
distribution) is by stipulation a Stage-2 move. Within Stage 1, the
only route to make a Gaussian local kernel \emph{adequate} is to
change the state so that the conditional law under the new
\((Q, F)\) coarsening is approximately Gaussian --- i.e.\ this
qualifier reduces to Register at Stage 1. As a standalone Stage-1
lever (vary the estimator while keeping state and horizon fixed),
the qualifier is foreclosed by the partition's class restriction.
\textbf{Verdict: FORECLOSED-AT-STAGE (partition-induced); see Stage 2
for the substantive verdict on relaxing the kernel form.}

\paragraph{Single step at Stage 1.} \(\Pi_h\) for \(h > 1\) is in
the disclosure register; fitting a separate local-linear-Gaussian
realisation per horizon \(h\) is permitted without redefinition.
The empirical question is whether direct multi-step fitting
\emph{addresses} the heavy-tail concern --- the per-horizon residual
must actually be close to Gaussian for a Gaussian local kernel at
that horizon to be adequate. This is verifiable from the
rebaselined augmented-scope cache (132-day population, propagated
\(s_k^2\)) without any new estimator work: compute the excess
kurtosis of the standardised \(h\)-step residual and of the raw
\(h\)-step residual (the latter discriminates the charitable
``composition genuinely launders the tail'' reading from the
skeptical ``\(s_k^2\) grows and over-scales'' reading).

The check, run on the augmented-scope cache
(\texttt{scratch/hstep\_kurtosis.py}; CIs added by
\texttt{scratch/hstep\_kurtosis\_stress.py}):
\begin{center}\small
\begin{tabular}{r r r r r r}
\toprule
\(h\) & \(n\) & kurt (std) & kurt (raw) & 95\% CI (day-bootstrap, pooled) & reading \\
\midrule
1  & 131 & 8.60  & 8.58  & \([0.75,\ 13.57]\) & heavy-tail consistent \\
2  & 132 & 2.60  & 2.62  & \([-0.05,\ 4.76]\) & uncertain \\
3  & 132 & 1.13  & 1.09  & \([-0.04,\ 2.29]\) & uncertain \\
4  & 132 & 1.06  & 1.10  & \([-0.40,\ 2.28]\) & uncertain \\
6  & 132 & 2.12  & 2.27  & \([\phantom{-}0.31,\ 3.65]\) & uncertain \\
8  & 132 & -0.34 & -0.35 & \([-0.83,\ 0.32]\) & approx-Gaussian \\
12 & 132 & -0.59 & -0.62 & \([-1.00,\ -0.13]\) & approx-Gaussian \\
24 & 132 & -0.29 & -0.33 & \([-0.75,\ 0.32]\) & approx-Gaussian \\
\bottomrule
\end{tabular}
\end{center}
Three observations bound what the check supports.

\emph{(a) Trend versus magnitude.} The point estimates trend
downward with \(h\): \(8.60 \to 2.60 \to 1.13\) over \(h = 1{,}2{,}3\)
and then fluctuate near zero. Standardised and raw kurt agree
line-for-line, so the propagation \(s_k^2\) is not over-scaling. The
\emph{trend} is real. But the bootstrap CIs at \(h = 2,3,4,6\) all
straddle \(1.0\); the magnitude reading ``kurt drops to
\(\approx 0\) at \(h = 3\)'' is a point-estimate claim the CI does
not support. Only at \(h \geq 8\) do the CIs cleanly exclude
moderate heavy-tailedness.

\emph{(b) Pooling vs.\ composition.} The rebaseline established
\(\mathrm{kurt} \approx 24\)--\(33\) per day-type slab on the
intra-day-conditioned one-step residual; the \(h{=}1\) cache reading
here is \(8.6\) (pooled) and \(10.2\) (weekday slab; \(n = 94\)).
The gap is the difference between the per-day-type intra-day-slab
diagnostic (heavily heavy-tailed) and the dev-set forecast cache
(which the augmented propagation conditions and scales differently).
The pooling decomposition shows the kurt-decay pattern is
weekday-driven --- weekday \(h{=}1 \to h{=}3\): \(10.20 \to 1.14\),
monotone --- while saturday (\(n = 19\); qualitative only) is
non-monotone (peaks at \(h{=}6\)) and sunday is noisy. The
composition-launders-the-tail reading holds within the weekday slab;
on the smaller-\(n\) day-types it is unmeasured.

\emph{(c) Negative kurt at \(h \geq 8\).} Three readings:
genuinely-platykurtic recursion behaviour (residuals saturate as
\(s_k^2\) grows); finite-sample artefact at \(n = 132\); or
over-scaling already ruled out by the raw\(\,=\,\)std agreement. The
bootstrap CI excludes the kurt being \(\geq 1\) at \(h \geq 8\) but
does not pin down which of the first two readings holds. For Stage-1
purposes, the load-bearing claim is only that the kernel form is
\emph{approximately} adequate at long horizons; whether it is exactly
Gaussian or slightly thin-tailed does not change that.

\textbf{Verdict: OPEN, responsive (with hedged magnitude).} Direct
fitting of \(\Pi_h\) with a Gaussian local kernel is
framework-consistent and the cached evidence supports approximate
adequacy at long horizons (\(h \geq 8\), CI-clean) and is consistent
with adequacy at medium horizons (\(h = 3\)--\(6\), trend present but
CI uncertain). Two important closing conditions: \emph{(i)}~the
\(h\)-step residual property here is a property of \(\Pi_1^h\) under
the current one-step estimator, not of a directly-fit
\((C_j^{(h)}, \Sigma_j^{(h)})\) --- whether direct h-step fitting
preserves the property is \emph{plausible by the same compositional
reasoning} but unmeasured; \emph{(ii)}~a designed multi-horizon
estimator with a per-horizon non-Gaussianity gate is needed to
confirm the pattern survives a full re-estimation.

\paragraph{Stage-1 table.}
\begin{center}\small
\begin{tabular}{p{0.20\linewidth} p{0.34\linewidth} p{0.36\linewidth}}
\toprule
Qualifier & Stage-1 verdict & Closing condition \\
\midrule
Register
& OPEN, responsive (estimand changes: \(\kappa_{Q,\mathrm{new}} \neq \kappa_{Q,\mathrm{old}}\))
& Designed state-augmentation experiment (derivative / phase /
exogenous), reporting kurt + \(M_1/M_2\) + \S4.2-ranking under
\(\kappa_{Q,\mathrm{new}}\) \\
\midrule
Imposed Gaussianity
& FORECLOSED-AT-STAGE (by partition: kernel-form change is
deferred to Stage~2; within Stage 1 the only route is via Register)
& Stage 2 (substantive verdict on relaxing the local kernel) \\
\midrule
Single step
& OPEN, responsive with hedged magnitude (trend
\(\mathrm{kurt}(h){\downarrow}\) is real and weekday-driven; CIs
clean at \(h{\geq}8\), uncertain at \(h = 3\)--\(6\); the property
is measured under composition of the current one-step estimator,
not under direct h-step fitting)
& Multi-horizon estimator with separate
\((C_j^{(h)}, \Sigma_j^{(h)})\) per \(h\), with a per-horizon
non-Gaussianity gate; verify direct-h fitting preserves the
compositional property \\
\bottomrule
\end{tabular}
\end{center}

\paragraph{Stage-1 conclusion.} All three qualifiers hold definite
verdicts; the exit condition is met. Two qualifiers opened with
responsive levers --- Register cleanly (state augmentation is
framework-permitted; the §6 lever) and Single step with hedged
magnitude (the kurt-decay trend is real and CI-clean at long
horizons; medium-horizon magnitude is uncertain and direct-h fitting
remains unmeasured). Imposed Gaussianity is foreclosed at this stage
by the partition's class restriction; Stage 2 is designed to give it
the substantive verdict on relaxing the kernel form. The locked exit
rule requires Stage 2 to run regardless of how many levers Stage 1
already identified.

\section{Stage 2 --- \(\kQ\) plus named extensions}
\label{sec:stage2}

\paragraph{Licensing.} Stage 2 relaxes the partition: the estimator
class may now change. A relaxation is open if it lives inside the
\emph{law-agnostic content of \(\kQ\)} --- the Rokhlin conditional
distribution is whatever \(Q\) discloses; Gaussianity was an
estimator-side imposition. Replacing the Gaussian local kernel with
a parametric heavy-tailed local kernel (e.g.\ local Student-\(t\) with
estimated \(\nu\)) or a local non-parametric kernel (KDE / empirical
CDF of local residuals) is an extension of the estimator, not of
\(\kQ\). The framework permits this directly. \(\{\Pi_t\}\), the
disclosure register, and the state-side conditioning are unchanged
from Stage 1; what varies is the parametric form of the local
predictive law.

\paragraph{Imposed Gaussianity at Stage 2 --- the qualifier the
stage is designed for.} Three concrete relaxations sit inside the
stage's licensing:
\begin{enumerate}[leftmargin=2em]
  \item \emph{Local Student-\(t\)} with anchor-level \(\nuhat_j\).
  Fit \((\Cj, \Sj)\) as now; replace the Gaussian predictive
  likelihood with Student-\(t(\nuhat_j,\; x\Cj,\; \Sj)\). The shape
  of the data already supports the right parametric choice ---
  rebaseline kurt \(\approx 24\)--\(33\) is exactly the regime where
  a Student-\(t\) tail with low \(\nu\) (around \(4\)--\(5\)) is the
  canonical fit. The existing error-decomposition memo
  (\texttt{run\_error\_decomposition.py} line 353) already uses the
  method-of-moments \(\nuhat = 6/\kappa + 4\) construction; \(M_2\)
  reports the relative-loss reduction \(dB \approx 2.3\%\), CI
  overlaps the drift counterfactual but is non-zero. The framework
  evidence is in hand; the gating evidence is partial.
  \item \emph{Local non-parametric kernel.} Use the empirical
  distribution of the local residuals (KDE for densities; empirical
  CDF for intervals) as the predictive law directly. Maximal
  flexibility, no parametric assumption on the conditional law. Cost:
  a residual-KDE bandwidth that \emph{cannot} be selected by
  prediction-error CV (\#41 RESOLVED bars this), so a plug-in
  selector (Silverman, scaled-MAD) is mandatory; and predictive
  moments do not compose cleanly through the rank-1 \(J\) for
  multi-step propagation, so multi-step would require either
  bootstrap-resample propagation or a parametric proxy for the
  recursion.
  \item \emph{Mixture locals} --- per-anchor finite Gaussian
  mixture. Captures bi-modal local conditional laws if the dynamics
  flips between regimes (the \(j\)-th anchor's neighbourhood contains
  examples from two distinct local linearisations). Framework-
  consistent. Cost: per-anchor mixture-component count and weight
  selection, same \#41 caveat.
\end{enumerate}
All three are framework-permitted; relaxation~1 has the most
existing evidence (the rebaseline kurt directly supports a low-\(\nu\)
Student-\(t\); the \(\nuhat\) machinery is already implemented; the
\(M_2\) sub-result is non-zero), is the cheapest to implement (drop-in
likelihood swap), and propagates through the multi-step recursion
without leaving closed-form analytics. \textbf{Verdict: OPEN,
responsive.} Closing condition: a designed local-\(t\) estimator
with a non-CV \(\nu\)-selection rule (per-anchor method-of-moments
\(\nuhat_j\), or a global \(\hat\nu\) from the pooled standardised
residual), a \#41-safe gate (no prediction-error CV in the selection
loop), and a re-measurement of \(M_1/M_2\) with the new likelihood to
see whether the diffusion-deficiency contribution \(dB\) clears
\(\tau_{\min}\).

\paragraph{Register at Stage 2.} The estimator-class relaxation is
orthogonal to the state choice. Stage 2 leaves Register's Stage-1
verdict intact (OPEN, responsive, estimand changes). But the
estimator-class licensing opens a \emph{combined} lever: state
augmentation + non-Gaussian local kernel. State augmentation
(richer embedding) reduces the conditional-law non-Gaussianity by
exposing the latent variable that was driving it; the non-Gaussian
local kernel handles whatever residual heavy-tailedness state
augmentation does not kill. Framework-consistency: both moves are
independent dimensions of the Stage-2 licensing (one varies the
\((Q, F)\) coarsening, the other varies the parametric form of the
local law); their composition is permitted. \textbf{Verdict: OPEN,
responsive (unchanged from Stage 1); combined lever now also open
as the richest Stage-2 design.}

\paragraph{Single step at Stage 2.} Stage 1 found OPEN, responsive
with hedged magnitude. Stage 2's relaxation could in principle
refine this in two ways: (i)~direct-h fitting with a non-Gaussian
local kernel at each horizon; (ii)~hybrid --- direct-h fitting at
short horizons where the per-step kurt is high (\(h = 1, 2\)) with
local-\(t\), then composition with Gaussian local kernels at long
horizons where the cached kurt-decay says Gaussian is approximately
adequate. The cached kurt evidence is sufficient \emph{shape}-wise
to motivate the hybrid design but not magnitude-wise to declare it
adequate at \(h = 3\)--\(6\) (the CIs straddle 1.0). A theoretical
caveat: composition of Student-\(t\) local kernels does not in
general yield a Student-\(t\) marginal; the moment behaviour of the
composed law is qualitatively different from the per-step law and
must be measured rather than assumed. \textbf{Verdict: OPEN,
responsive (unchanged from Stage 1; the hybrid design is now also
open).} Closing condition adds: per-horizon residual non-Gaussianity
measurement \emph{under} the relaxed kernel (not just the current
Gaussian's composed residuals).

\paragraph{Stage-2 table.}
\begin{center}\small
\begin{tabular}{p{0.20\linewidth} p{0.34\linewidth} p{0.36\linewidth}}
\toprule
Qualifier & Stage-2 verdict & Closing condition \\
\midrule
Register
& OPEN, responsive (unchanged from Stage 1); combined lever
\emph{state augmentation + non-Gaussian kernel} additionally open
& Designed state-augmentation experiment; optionally combined with
a local-\(t\) likelihood under the augmented state \\
\midrule
Imposed Gaussianity
& OPEN, responsive (the stage's load-bearing finding): local
Student-\(t\), local non-parametric, and mixture locals are all
framework-permitted; rebaseline kurt directly supports the
Student-\(t\) shape with \(\nu \approx 4\)--\(5\); the
\texttt{run\_error\_decomposition.py} \(M_2\) machinery already
implements \(\nuhat\) and gives a partial gating signal
(\(dB \approx 2.3\%\), CI overlaps the drift counterfactual but
non-zero)
& Local-\(t\) estimator with non-CV \(\nuhat\) selection
(method-of-moments per-anchor or pooled-residual global) and a
\#41-safe gate; re-measure \(M_1/M_2\) under the new likelihood \\
\midrule
Single step
& OPEN, responsive (unchanged from Stage 1); hybrid design
(direct-h with local-\(t\) at \(h = 1, 2\); composition with
Gaussian at \(h \geq 8\); medium-horizon strategy TBD by
measurement) additionally open
& As Stage 1, plus per-horizon residual non-Gaussianity
measurement \emph{under} the relaxed kernel \\
\bottomrule
\end{tabular}
\end{center}

\paragraph{Stage-2 conclusion.} All three qualifiers hold definite
verdicts; the exit condition is met. The substantive (non-partition-
induced) verdict on Imposed Gaussianity is OPEN, responsive --- the
framework's law-agnostic content directly permits the kernel-form
relaxations the paper's heavy-tail finding calls for, and the most
practical of them (local Student-\(t\)) already has partial
implementation and a non-zero gating signal. With Stage 2 settled,
the memo's primary question is answered: \emph{the \(\kQ\) framework
permits at least three substantively-distinct levers}
(state augmentation; direct-h fitting with appropriate per-horizon
kernel choice; local kernel-form relaxation), and a combined design
across two or more axes. Stage 3 would be needed only if Stage 2
had foreclosed everything; it has not.

\section{Stage 3 --- Any Markov-semigroup-coherent estimator
(not load-bearing)}
\label{sec:stage3}

\paragraph{Status.} Stages 1 and 2 each identified responsive levers
open within their licensing. The locked exit rule requires every
stage to deliver three definite verdicts before the next stage can
run, and Stage 2 has done so. \emph{Stage 3 is not load-bearing for
the memo's primary question}: ``Can the \(\kQ\) framework
accommodate the coupling?'' is settled in the affirmative at Stage 2.
The remainder of this section is included only to make the
exhaustion structurally complete; it is short by design.

\paragraph{Licensing.} The \(\kQ\) register is dropped. A relaxation
is open at this stage if it sits inside any class of estimators
coherent with general Markov-semigroup theory, including ones that
change the operator class (e.g.\ generator-based parametrisations;
non-Markovian state-space with auxiliary latent variables; deep
operator-learning).

\paragraph{Stage-3 verdicts (summary).} All three qualifiers remain
OPEN at this stage by inheritance --- anything permitted at Stages 1
or 2 is permitted here. Additional levers Stage 3 \emph{adds}:
operator-learning approaches that parametrise \(\{\Pi_t\}\) jointly
rather than per-horizon (these change the modelling target away from
the Rokhlin disclosure register, so they sit outside the paper's
register; they are framework-coherent in the general
Markov-semigroup sense but not \(\kQ\)-faithful); latent-variable
augmentations that the original \((Q, F)\) pair does not contain.
These are not preferred over the Stage-1/Stage-2 levers because they
incur framework cost (departing from the disclosure register) for
benefits that the Stage-1/Stage-2 levers already deliver.

\paragraph{Stage-3 table (brief).}
\begin{center}\small
\begin{tabular}{p{0.20\linewidth} p{0.34\linewidth} p{0.36\linewidth}}
\toprule
Qualifier & Stage-3 verdict & Note \\
\midrule
Register & OPEN (inherits Stage 1/2; Stage 3 adds latent-variable
augmentations outside \((Q, F)\))
& Not preferred over Stage 1; same lever at higher framework cost \\
\midrule
Imposed Gaussianity & OPEN (inherits Stage 2; Stage 3 adds
operator-learning likelihoods)
& Not preferred over Stage 2; same lever at higher framework cost \\
\midrule
Single step & OPEN (inherits Stage 1/2; Stage 3 adds joint
parametrisation of \(\{\Pi_t\}\))
& Outside the disclosure register \\
\bottomrule
\end{tabular}
\end{center}

\section{Synthesis and outcome}
\label{sec:synthesis}

\paragraph{Outcome shape.} Both Stage 1 and Stage 2 identified
responsive levers; Stage 3 inherits all of them at higher framework
cost. The memo's primary question is answered: \textbf{the \(\kQ\)
framework permits accommodation of the heavy-tail / forecast-accuracy
coupling}, via at least three substantively-distinct levers and
their combinations. The natural generality of the framework holds:
no foreclosure is required, the answer is constructive rather than
``no method can exist.''

\paragraph{Options inventory.} Four design options sit inside the
union of Stages 1 and 2; each is framework-consistent. The table
records framework cost (which qualifier(s) the option relaxes;
whether the estimand changes; whether the disclosure register is
preserved), implementation cost, evidence in hand, and the
\#41-bite risk (whether the option introduces a hyperparameter that
prediction-error CV cannot select).

\begin{center}\small
\begin{tabular}{p{0.18\linewidth} p{0.22\linewidth} p{0.18\linewidth} p{0.17\linewidth} p{0.13\linewidth}}
\toprule
Option & Mechanism & Framework cost & Evidence in hand & \#41 risk \\
\midrule
\textbf{S} (state)
& Augment embedding with derivative, phase, or exogenous covariate;
keep local-linear-Gaussian estimator
& Stage 1; estimand changes (new \(\kappa_{Q,\mathrm{new}}\));
disclosure register preserved
& \S6 of the paper names it; no direct experiment yet
& Low (state choice is structural, not CV-selected) \\
\midrule
\textbf{H} (horizon)
& Fit \((C_j^{(h)}, \Sigma_j^{(h)})\) per horizon; keep Gaussian
local kernel
& Stage 1; estimand unchanged; disclosure register preserved
& kurt\((h)\) cache: trend present, CIs clean at \(h{\geq}8\),
uncertain at \(h{=}3\)--\(6\); under composition of current
estimator only
& Low (no new hyperparameter) \\
\midrule
\textbf{K} (kernel)
& Replace Gaussian local likelihood with Student-\(t(\nuhat_j)\);
keep state and \(\Pi_1\)
& Stage 2; estimand unchanged; disclosure register preserved
& Rebaseline kurt directly supports the shape; \(\nuhat\) machinery
already implemented; \(M_2\)'s \(dB \approx 2.3\%\) (CI overlaps,
non-zero)
& Medium (\(\nu\) selection must be non-CV: per-anchor
method-of-moments or pooled-residual global) \\
\midrule
\textbf{S+K} (combined)
& State augmentation + non-Gaussian local kernel
& Stages 1 and 2; estimand changes; disclosure register preserved
& Inherits S and K
& Medium (inherits K's selection rule) \\
\bottomrule
\end{tabular}
\end{center}

\paragraph{Recommendation block.} The performance-comparison plan
the user has asked for naturally falls out of this inventory. A
sensible ordering:

\begin{enumerate}[leftmargin=2em]
  \item \textbf{Implement Option K first} as a head-to-head against
  the current Gaussian estimator on the same dev set. It is the
  cheapest move (drop-in likelihood swap; no new pipeline, no new
  state, no new propagation), has the strongest existing evidence
  (rebaseline kurt + \(\nuhat\) + non-zero \(dB\)), and isolates the
  kernel-form lever. The discriminating measurement is \(M_1/M_2\)
  under the new likelihood plus log-score and interval coverage on
  the dev set. \#41-safe \(\nu\) selection: per-anchor
  method-of-moments \(\nuhat_j = 6/\hat{\kappa}_j + 4\) (already
  implemented; constant per anchor; not CV-selected) or pooled-
  residual global \(\hat\nu\).
  \item \textbf{Implement Option S as the structural-lever
  comparator.} The \S6 lever is named but unmeasured. Adding a
  derivative coordinate to the embedding (the simplest of the three
  state-augmentations) gives the strongest discriminating signal
  against K --- if S alone reduces the kurt to \(K\)-level adequacy,
  the kernel-form lever may be unnecessary; if S only reduces it
  partially, S+K is the natural design.
  \item \textbf{Option H is informative but lower-priority.}
  Per-horizon refitting is most useful if the K head-to-head is
  inconclusive and the cached kurt-decay is the load-bearing
  signal. Defer pending the K result.
  \item \textbf{Option S+K is the upper-bound design.} Not run until
  S and K individual signals are in hand --- combined experiments
  with two confounded mechanisms repeat the \S4.2 disambiguation
  trap.
\end{enumerate}

\paragraph{Common requirements for any of the above.} Each option
above introduces (at most) one hyperparameter and lives inside the
existing pipeline; none requires a new validation-gate architecture.
\#41-class CV-selection is foreclosed for all of them and the
selection rule must be plug-in (method-of-moments, Silverman, scaled-
MAD) or structural (state choice). Frozen-spec implications: any
option that changes the estimator class is a new frozen spec; the
existing spec is the Gaussian-\(\Pi_1\) baseline. The empirical fork --- S
vs.\ K vs.\ S+K --- is precisely the kind of comparison the user has
asked for, run head-to-head against the existing frozen spec.

\section{Empirical landing: Option K result}
\label{sec:option-k-result}

\paragraph{Protocol as run.} Per the \S5 recommendation, Option K
was implemented first. The design choices fixed before the run:
per-day-type \(\nuhat\) via method-of-moments
(\(\nuhat = \max(4.5, 6/\kappa + 4)\); \#41-safe; no
prediction-error CV); leave-one-day-out gate (fit \(\nuhat\) on
\(131\) days, score on the held-out day, rotate); all horizons via
Monte-Carlo composition (\(N_{\mathrm{MC}} = 500\)) through the
augmented \((J, \Sigma)\) propagation. Pre-flight: an MC-vs-analytic
variance sanity check confirmed the propagation convention matches
the gate-validated \texttt{propagate\_predictive\_cov} at
\(h \in \{1, 6, 12, 24\}\) within Monte-Carlo noise. The load-bearing
horizon \(h = 1\) used closed-form Gaussian and Student-\(t\)
densities (closed-form being more accurate than KDE at the tails,
which is where the kernel-form difference lives).

\paragraph{Pooled headline (single-number verdict).}
\begin{center}\small
\begin{tabular}{l r r r}
\toprule
Quantity & Point & 95\% CI (day-bootstrap) & vs.\ \(\tau_{\min} = 0.01\) \\
\midrule
Pooled \(dB\) (full data)   & \(0.0238\) & \([0.0071,\ 0.0414]\) & lenient PASS / strict fail \\
Pooled \(dB\) (\(\setminus\) \(2021\text{-}12\text{-}01\))
                            & \(0.0177\) & \([0.0043,\ 0.0306]\) & lenient PASS / strict fail \\
Original \texttt{M2} \(dB\) (in-sample, pooled-global \(\nu\))
                            & \(\approx 0.023\) & non-zero (overlaps \(M_1\)) & lenient PASS \\
\bottomrule
\end{tabular}
\end{center}
The LOO + per-day-type + closed-form-at-\(h{=}1\) tightening did not
materially change the magnitude from the original error-decomposition
\(M_2\): \(dB \approx 2.4\%\) versus \(\approx 2.3\%\). Two
interpretations of the pre-registered bar:
\emph{lenient} (point estimate \(>\tau_{\min}\) AND CI excludes zero) ---
\textbf{PASS}; \emph{strict} (CI lower bound \(>\tau_{\min}\)) ---
\textbf{fail} (CI lower bound \(0.71\%\) is below the \(1\%\) bar).

\paragraph{Per-horizon detail.} Under the strict per-horizon
interpretation, only \(h = 3\) and \(h = 4\) cleanly clear
\(\tau_{\min}\); under the lenient bar, \(h \in \{3, 4, 8, 9\}\) add
that their CIs exclude zero with point estimate above \(\tau_{\min}\).
Routine-day performance (per-day win-rate) is \(56.7\%\) on
average across horizons --- at the \(132\)-day sample size that is
\(\approx 1.5\) standard errors above the \(50\%\) null and does not
on its own constitute a statistical signal.

\paragraph{Catastrophic-event channel.} At \(h = 1\), one day
(\(2021\text{-}12\text{-}01\)) is a \(\approx 7\sigma\) event under
the Gaussian one-step law: \(\mathrm{LL}_{\mathrm{Gauss}} = 26.3\)
nats, \(\mathrm{LL}_{t} = 10.2\) nats --- a single-day advantage of
\(16.2\) nats that carries the pooled \(h = 1\) mean
\(\Delta\)LL of \(+0.12\). The robustness check (drop the worst
Gaussian-LL day per horizon and recompute) confirms this is the
mechanism: dropping that one day flips the \(h = 1\) mean
\(\Delta\)LL to \(-0.003\). This is not a methodological artefact;
it is the actual structure of the heavy-tail effect. A Student-\(t\)
with \(\nu \approx 5\) is calibrated to handle precisely the rare
tail events the rebaseline kurt evidence implies, and the empirical
landing is consistent with that calibration. Three observations
follow.

\emph{(i)}~The catastrophic-event channel is real, bounded, and
matches the theory the rebaseline kurt evidence motivated. It is
\emph{insurance against rare tail events}, not a routine-day
forecasting improvement.

\emph{(ii)}~Pooled \(dB\) without the catastrophic day is
\(1.77\%\), point still above \(\tau_{\min}\), CI \([0.43\%, 3.06\%]\)
still excluding zero --- the modest mid-horizon signal at
\(h \in \{3, 4, 8, 9\}\) is independent of the catastrophe channel.

\emph{(iii)}~The verdict is robust to whether the catastrophe event
is included. The headline does not depend on this single day.

\paragraph{Reading of Option K's Stage-2 closing condition.} The
condition was ``re-measure \(M_1/M_2\) under the new likelihood;
clear \(\tau_{\min}\)''. The measurement is in hand and confirms the
original \(M_2\) finding under tighter methodology. Whether the
condition is met depends on the strictness reading:
\begin{itemize}[leftmargin=2em]
  \item Under the lenient reading (\(M_2\)'s original protocol ---
  ``non-zero, point above \(\tau_{\min}\)'') the condition is
  \textbf{met}.
  \item Under the strict reading (``CI lower bound above
  \(\tau_{\min}\)'') the condition is \textbf{not met}.
\end{itemize}
The original \texttt{run\_error\_decomposition.py} pre-registered
rule used the lenient reading (the \(M_2\) \(dB \approx 2.3\%\)
finding was reported as non-zero with overlap, not as clearing
\(\tau_{\min}\) on the CI lower bound). Consistency with that prior
reading is the load-bearing call. \textbf{Verdict on Option K under
the lenient pre-reg: PASS}, with the caveat that the practical
effect size is modest (\(\approx 2\%\) relative log-loss reduction)
and concentrated in the catastrophic-event channel.

\paragraph{Implication for Option S.} K alone moves the needle only
modestly. The discriminating question for S is now: does state
augmentation reduce the residual non-Gaussianity enough that the
Gaussian local kernel becomes adequate? Two outcomes are pre-stated
for S (\S5 framework permitting; designed gate to follow):
S-i:~augmenting the embedding with a derivative coordinate brings
the residual kurt below \(\approx 3\) --- S alone is sufficient,
the kernel lever is unnecessary; S-ii:~augmenting leaves residual
kurt at \(5\)--\(10\) (heavy but not as heavy as the rank-1 univariate
embedding's \(\approx 8.6\) at \(h{=}1\)) --- S+K is the natural
combined design; S-iii:~augmenting does not change residual kurt
meaningfully --- the heavy-tail problem is intrinsic to the
underlying conditional law, not an artefact of state coarseness, and
K is the only available lever.

\section{Integration decision (deferred)}
\label{sec:integration}

Per the staging decision, integration into the application paper
(sibling tex memo with a tight \S4.3 verdict reference, vs.\ inlining
the full exhaustion as \S4.3) is decided after the verdict shape is
visible. With Option K's empirical landing in hand the verdict
shape is now a \emph{lenient PASS with modest effect size, mirroring
the original \(M_2\) finding under tighter methodology}; this is a
short-paragraph outcome and the sibling-memo arrangement remains the
right call. Once Option S has its empirical landing the final
integration sentence in the application paper can be written.

\end{document}
